% Options for packages loaded elsewhere
\PassOptionsToPackage{unicode}{hyperref}
\PassOptionsToPackage{hyphens}{url}
%
\documentclass[
]{article}
\usepackage{amsmath,amssymb}
\usepackage{lmodern}
\usepackage{iftex}
\ifPDFTeX
  \usepackage[T1]{fontenc}
  \usepackage[utf8]{inputenc}
  \usepackage{textcomp} % provide euro and other symbols
\else % if luatex or xetex
  \usepackage{unicode-math}
  \defaultfontfeatures{Scale=MatchLowercase}
  \defaultfontfeatures[\rmfamily]{Ligatures=TeX,Scale=1}
\fi
% Use upquote if available, for straight quotes in verbatim environments
\IfFileExists{upquote.sty}{\usepackage{upquote}}{}
\IfFileExists{microtype.sty}{% use microtype if available
  \usepackage[]{microtype}
  \UseMicrotypeSet[protrusion]{basicmath} % disable protrusion for tt fonts
}{}
\makeatletter
\@ifundefined{KOMAClassName}{% if non-KOMA class
  \IfFileExists{parskip.sty}{%
    \usepackage{parskip}
  }{% else
    \setlength{\parindent}{0pt}
    \setlength{\parskip}{6pt plus 2pt minus 1pt}}
}{% if KOMA class
  \KOMAoptions{parskip=half}}
\makeatother
\usepackage{xcolor}
\setlength{\emergencystretch}{3em} % prevent overfull lines
\providecommand{\tightlist}{%
  \setlength{\itemsep}{0pt}\setlength{\parskip}{0pt}}
\setcounter{secnumdepth}{-\maxdimen} % remove section numbering
\ifLuaTeX
  \usepackage{selnolig}  % disable illegal ligatures
\fi
\IfFileExists{bookmark.sty}{\usepackage{bookmark}}{\usepackage{hyperref}}
\IfFileExists{xurl.sty}{\usepackage{xurl}}{} % add URL line breaks if available
\urlstyle{same} % disable monospaced font for URLs
\hypersetup{
  hidelinks,
  pdfcreator={LaTeX via pandoc}}

\author{Dietmar Gerald Schrausser}
\date{09.05.2026}


\begin{document}

\hypertarget{foliumblat-viv}{%
\section{Foliu{[}m{]}/Blat VIv}\label{foliumblat-viv}}

Eine zeitgenössische deutsche Übersetzung wird präsentiert, die auf
einem Vergleich der bestehenden lateinischen, frühneuhochdeutschen
(ENHG) und englischen (Übersetzungs-) Texte basiert, um ein besseres
Verständnis zu ermöglichen. Insbesondere da der ENHG-Text im Vergleich
zum modernen Hochdeutschen schwer verständlich ist, da dieser fundierte
Deutschkenntnisse (als Muttersprache) sowie Kenntnisse verschiedener
Dialekte voraussetzt.

\hypertarget{das-erste-zeitalter-der-welt}{%
\subsection{Das erste Zeitalter der
Welt}\label{das-erste-zeitalter-der-welt}}

Das erste Zeitalter der Welt von \emph{Adam} bis zur \emph{Sintflut}
umfasste laut den Hebräern 1656 Jahre, gemäss der \textbf{Septuaginta},
\textbf{Isidor} und vielen anderen, waren es hingegen 2242 Jahre, was
folglich auch als gängige Zeitrechnung übernommen wurde.

Die \emph{höchste Güte}\footnote{`SUmma bonitas' und `DIe höhst
  gůthait'.} will diese Güte auch anderen mitteilen, darum erschuf sie
ein vernunftbegabtes Wesen, welches die höchste Güte verstünde und
dadurch liebte\footnote{`i{[}n{]}tellige{[}n{]}do amaret' und
  `versteende liebhet'.}, {[}die höchste Güte{]} durch diese Liebe
besäße und gesegnet\footnote{`beata' und `selig'.} sei.\footnote{vom
  `Sententiarum' (Petrus,
  \href{https://doi.org/10.3931/e-rara-18241}{1484}, Buch II, Kapitel
  4). Diese Formulierung bezieht sich auf eine theologische Aussage, die
  oft auch dem heiligen Augustinus zugeschrieben wird und den Zweck der
  \emph{vernünftigen Seele} erklärt: `Fecit \emph{Deus} rationalem
  creaturam ut summum bonum intelligeret, intelligendo amaret, amando
  possideret, possidendo frueretur.'}

Gott schuf \emph{nun}\footnote{`aut{[}em{]}' und `Aber'.} den ersten
Menschen, indem er seinen Leib mit Hilfe der Engel aus dem Lehm der Erde
auf dem Feld von \emph{Damaskus} formte und sein \emph{Gesicht} mit dem
\emph{Atem des Lebens}\footnote{`spiraculum vite' und `geystung des
  lebens'.} inspirierte\footnote{`inspirauit' und `eingeystet'.}, das
heißt\footnote{charakteristisch für das `Compendium historiae in
  genealogia Christi' (Poitiers,
  \href{https://libwww.freelibrary.org/digital/item/22446}{1499}).}, er
schuf eine Seele\footnote{`a{[}n{]}i{[}m{]}am' und `sel'.}, die er mit
dem geformten Leib vereinigte. So wurde der Mensch in seinen natürlichen
Eigenschaften nach dem \emph{Bild}\footnote{`ymaginem' und `pildnus'.}
Gottes und in seinen geistigen Gaben\footnote{`gratuitis' und
  `gnadenreichen dingen'.} nach seinem \emph{Gleichnis}\footnote{`si{[}mi{]}litudinem'
  und `gleichnuß'.} geschaffen. Siehe, der Herr lässt uns wundervolles
Wohlwollen\footnote{`gra{[}tia{]}' und `gnad'.} zukommen.

\begin{quote}
Cu{[}m{]} ergo mare{[}m{]} ad si{[}mi{]}litudine{[}m{]} sua{[}m{]}
primum finxisset: tu{[}m{]} etia{[}m{]} feminam {[}con{]}figurauit
adipsius hominis effigiem vt dou inter se p{[}er{]}mixti sex{[}us{]}
p{[}ro{]}pagare sobolem possent: {[}et{]} omne{[}m{]} terra{[}m{]}
multitudine opplere.\footnote{`Als er also zuerst den Mann nach seinem
  Ebenbild geschaffen hatte, schuf er auch die Frau nach dem Bild des
  Mannes selbst, damit die beiden Geschlechter, vereint, Nachkommen
  zeugen und die ganze Erde mit einer Vielzahl füllen könnten', aus den
  `Divine Institutes' (Lactantius,
  \href{https://books.google.com/books?id=qGsnghsi3uUC}{1497}, Buch 2,
  Kapitel 10).}
\end{quote}

\begin{quote}
FOrmatis a{[}n{]}i{[}m{]}antib{[}us{]} terre {[}et{]} volatilib{[}us{]}
adduct{[}is{]} ad Ada{[}m{]} vt videret ea: cu{[}m{]} ada{[}m{]}
no{[}n{]} i{[}n{]}ueniret{[}ur{]} adiutor si{[}mi{]}lis sibi: immisit
d{[}omi{]}n{[}u{]}s sopore{[}m{]} in Ada{[}m{]} {[}et{]} tulis
una{[}m{]} de costis ei{[}us{]}: reple{[}n{]}s carne{[}m{]} pro ea:
{[}et{]} edificauit in muliere{[}m{]}. Qua{[}m{]} ada{[}m{]}
vide{[}n{]}s: dixit h{[}oc{]} nunc os de ossib{[}us{]} meis h{[}oc{]}
vocabit{[}ur{]} Issa q{[}uo{]}d lati{[}n{]}e m{[}u{]}l{[}ie{]}r
interp{[}re{]}tat{[}ur{]}: q{[}uorum{]} de viro su{[}m{]}pta
e{[}st{]}.\footnote{`Nachdem Gott die Lebewesen der Erde und die Vögel
  erschaffen hatte, brachte er sie zu Adam, damit dieser sie sähe. Da
  sich für Adam keine Gehilfin fand, die ihm glich, ließ der Herr Adam
  in einen tiefen Schlaf fallen, nahm eine seiner Rippen, füllte sie mit
  Fleisch und formte daraus eine Frau. Als Adam sie sah, sprach er: Dies
  ist nun Bein von meinem Bein; sie soll \emph{Issa} heißen, was auf
  Lateinisch als Frau übersetzt wird, weil sie vom Mann genommen wurde',
  Hebräisch \emph{Ishah} (Frau) von \emph{Ish} (Mann), `Vulgata',
  Genesis 2,19-23 (vgl. Jiménez de Cisneros,
  \href{https://doi.org/10.3931/e-rara-46695}{1517}).}
\end{quote}

Nachdem Gott Adam erschaffen hatte, führte er ihn ins Paradies und
machte dort aus der \emph{Rippe}\footnote{`costa' und `ripp', Hebräisch
  \emph{tsela} (Seite).} des Schlafenden \emph{Eva}, formte ihm so eine
Gefährtin. Er machte sie weder vom \emph{Kopf}\footnote{`capite' und
  `haubt'.} her, so dass sie nicht über den Mann herrsche, noch vom
\emph{Fuß} des Mannes, damit sie nicht verschmäht werde, sondern von der
\emph{Seite}\footnote{`latere' und `seyten'.}, damit sich das Band der
Liebe bewährt\footnote{`p{[}ro{]}baret{[}ur{]}' und `zu bewerung'.}.\footnote{vom
  `Sententiarum' (Petrus,
  \href{https://doi.org/10.3931/e-rara-18241}{1484}, Buch IV), ein
  gängiges Motiv in der mittelalterlichen Scholastik, siehe Thomas von
  Aquin (\href{https://doi.org/10.3931/e-rara-15170}{1489}) oder Hugo
  von St.~Viktor
  (\href{https://archive.org/details/hugh-of-st.-victor-roy-j.-deferrari-on-the-sacraments-of-the-christian-faith-de-/mode/1up}{1951},
  Buch I, Teil VI, Kapitel XXXV).}

Damit wird auch gezeigt, dass nicht der Wert\footnote{`nobilitate' und
  `adel'.} des Herkunftsortes, sondern die eigene \emph{Güte}\footnote{`virtute'
  und `tugend'.} jedem einzelnen\footnote{`vnusq{[}ui{]}sq{[}ue{]}' und
  `ein yder'.} Menschen \emph{moralischen Charakter}\footnote{`gra{[}tia{]}m'
  und `gnad'.} verleiht. Der Mann, außerhalb des Paradieses erschaffen,
erweist sich, obwohl an einem \emph{minderen} Ort\footnote{`inferiori
  loco' und `vndern stat'.} geschaffen, dennoch \emph{besser
für}\footnote{`melior inuenit{[}ur{]} p{[}ro{]}' und `besser gefunden
  für'.} Eva, welche im Paradies erschaffen wurde.\footnote{aus `De
  Paradiso' (Ambrosius,
  \href{https://books.google.com/books?id=L1DbnAAACAAJ}{1897}, 4.24),
  ein gängiges theologisches Argument hinsichtlich des Verhältnisses
  zwischen dem Herkunftsort eines Menschen und seinem moralischen
  Charakter.}

Gott, der oberste Schöpfer aller Dinge, schuf somit am sechsten Tag der
Welt, am 25. März\footnote{in der mittelalterlichen Theologie
  traditionell sowohl als Tag der Weltschöpfung als auch als Tag der
  Kreuzigung angesehen.}, nach den Tieren der Erde, allen Kriechtieren
und Vögeln, \emph{Adam}, den zuerst geschaffenen\footnote{`p{[}ro{]}thoplastu{[}m{]}',
  ein latinisierter griechischer Begriff (protoplast).}, ersten
Menschen, aus einem roten Erdklumpen oder Lehm\footnote{vom hebräischen
  Wort \emph{adamah} (Erde/Boden).} auf dem Feld von
Damaskus\footnote{nach mittelalterlichen Legenden war die Erschaffung
  Adams auf einem bestimmten Feld in der Nähe von Damaskus oder Hebron.},
\emph{als Vollendung}\footnote{`fine{[}m{]}' und `ein end'.} \emph{und
Inhaber}\footnote{`possessore{[}m{]}' und `besitzer'.} aller Geschöpfe.

\hypertarget{references}{%
\subsection{References}\label{references}}

Ambrosius, A. (1897). \emph{Opera: Exameron. De Paradiso. De Cain Et
Abel. De Noe. De Abraham. De Isaac. De Bono Mortis}. Edited by Faller,
O., \& Schenkl, K. Corpus Scriptorum Ecclesiasticorum Latinorum.
Hoelder-Pichler-Tempsky.
\url{https://books.google.com/books?id=L1DbnAAACAAJ}

Hugh of St.~Victor. (1951). \emph{On the Sacraments of the Christian
Faith (de Sacramentis)}. Edited by Deferrari, R J. Cambridge,
Massachusetts: The Mediaeval Academy Of America.
\url{https://archive.org/details/hugh-of-st.-victor-roy-j.-deferrari-on-the-sacraments-of-the-christian-faith-de-/mode/1up}

Jiménez de Cisneros, F. (1517). \emph{Biblia Polyglotta Complutensis}.
Complutum: Arnaldo Guillén de Brocar.
\url{https://doi.org/10.3931/e-rara-46695}

Lactantius, L.C.F. (1497). \emph{Divinae Institutiones}. Venetiis: Simon
Bevilaqua. \url{https://books.google.com/books?id=qGsnghsi3uUC}

Petrus. (1484). \emph{Sententiarum Libri IV}. Basel.
\url{https://doi.org/10.3931/e-rara-18241}

Poitiers P. (1499). \emph{Compendium Historiae in Genealogia Christi}.
Philadelphia, PA: Free Library of Philadelphia.
\url{https://libwww.freelibrary.org/digital/item/22446}

Thomas. (1485). \emph{Summa Theologica: Partes 1-3}. Basel: Michael
Wenssler. \url{https://doi.org/10.3931/e-rara-15170}

\end{document}
